% Pot catalytique : monolithe en nid d'abeilles et étapes de la catalyse en surface (exercice 14, chapitre 5)
\begin{tikzpicture}[x=1mm,y=1mm,
  met/.style={circle,ball color=otInk!30,inner sep=0pt,minimum size=4.2mm},
  aC/.style={circle,fill=otInk!85,inner sep=0pt,minimum size=2.4mm},
  aO/.style={circle,fill=otRed!75,inner sep=0pt,minimum size=2.6mm},
  aN/.style={circle,fill=otNavy,inner sep=0pt,minimum size=2.6mm},
  mouv/.style={otGreenDark,line width=0.7pt,-{Stealth[length=1.8mm]}}]
% monolithe
\draw[otInk!70,line width=0.7pt,rounded corners=2mm,fill=otRule!25] (0,0) rectangle (34,26);
\foreach \i in {0,...,6}{\foreach \j in {0,...,4}{\filldraw[fill=white,draw=otInk!45,line width=0.35pt] (2.2+\i*4.4,2.4+\j*4.4) rectangle ++(3.4,3.4);}}
\draw[otMuted,line width=1pt,-{Stealth[length=2.4mm]}] (-13,13) -- (-2,13);
\node[pcLegende,anchor=south,align=center] at (-7.5,14) {gaz\\polluants};
\draw[otMuted,line width=1pt,-{Stealth[length=2.4mm]}] (36,13) -- (45,13);
\node[pcLegende,anchor=south,align=center] at (40.5,14) {gaz\\traités};
\node[pcLegende,anchor=north,align=center] at (17,-1.5) {céramique en nid d'abeilles\\recouverte de Pt, Pd et Rh};
% surface du métal
\begin{scope}[shift={(52,0)}]
\foreach \k in {0,...,19}{\node[met] at (\k*4.2,2.1) {};}
% 1 : adsorption de CO et de NO
\node[aC] at (5,19) {}; \node[aO] at (7.4,20.6) {};
\draw[mouv] (6,16.5) -- (6,9.6);
\node[aC] at (6,6.2) {}; \node[aO] at (6,8.5) {};
\node[aN] at (14,19.5) {}; \node[aO] at (16.4,21) {};
\draw[mouv] (15,17) -- (15,9.6);
\node[aN] at (15,6.2) {}; \node[aO] at (15,8.6) {};
\node[pcLegendeG] at (10.5,27) {1. adsorption};
% 2 : réaction en surface (O + CO, puis N + N)
\node[aO] at (27,6.2) {}; \node[aC] at (33,6.2) {}; \node[aO] at (33,8.5) {};
\draw[mouv] (27.3,8.6) .. controls (28,12.5) and (31.5,12.5) .. (32.2,10.2);
\node[aN] at (40,6.2) {}; \node[aN] at (46,6.2) {};
\draw[mouv] (40.3,8.6) .. controls (41,12) and (44.6,12) .. (45.4,8.8);
\node[pcLegendeG] at (36.5,27) {2. réaction en surface};
% 3 : désorption de CO2 et de N2
\node[aO] at (56,12) {}; \node[aC] at (58.5,12) {}; \node[aO] at (61,12) {};
\draw[mouv] (58.5,14.5) -- (58.5,20.5);
\node[pcLegende,anchor=west] at (62.6,12) {\ce{CO2}};
\node[aN] at (71,15) {}; \node[aN] at (73.5,15) {};
\draw[mouv] (72.2,17.5) -- (72.2,22.5);
\node[pcLegende,anchor=west] at (75,15) {\ce{N2}};
\node[pcLegendeG] at (65,27) {3. désorption};
\node[pcLegende,anchor=north] at (40,-1.5) {surface du métal (atomes de platine ou de rhodium)};
\end{scope}
\end{tikzpicture}
\par\smallskip\legende{Les molécules se fixent sur les atomes de métal (adsorption), y réagissent, puis les produits se détachent (désorption).}
